\maketitle
\vspace{-1.5em}

\begin{abstract}
Reliable reinforcement learning must work across initial states and independent training seeds. We evaluate five trained actors on the same 2,501-state Pendulum grid. The selected reinforcement-learning-plus-supervision recipe combines learner-state DAgger, automatic follow-up on measured weak starts, and conservative local critic search. It succeeds near the reference on 12,496 of 12,505 seed-state trials, or 99.928\%. The selected pure-RL recipe combines SimbaV2, FastSACN8 critic training, symmetry, and twin-critic action search. It reaches 12,008 trials, or 96.026\%. Task-success rates are 93.858\% and 92.875\%; strict reference wins are 12.555\% and 22.823\%. Matched ablations show that critic search repairs many decisions missed by the pure-RL actor, while the raw actor becomes less reliable and more saturated. On failure-conditioned states, multi-step actors retain about 7\% of a nominal local critic step after action-bound projection. The reference-prefix intervention repairs 19.1\% of failures with one state-matched action and 99.0\% with 32, showing that the selected raw-actor failures generally require an ordered, state-matched correction sequence.
\end{abstract}

\section{Reliability target and evidence}

\paragraph{Chasing two kinds of nines.}
The project asks how to increase reliability within one trained model and across random training seeds. A high pooled score can conceal different failures in different seeds. Every mainline method is therefore evaluated on 61 angles in $[-\pi,\pi]$ and 41 angular velocities in $[-1,1]$. Five actor seeds yield 12,505 deterministic 200-step rollouts in Gymnasium \texttt{Pendulum-v1}.

The observation, torque, and reward are
\[
s=(\cos\theta,\sin\theta,\dot\theta),\quad a\in[-2,2],
\]
\[
r_t=-(\theta_t^2+0.1\dot\theta_t^2+0.001a_t^2).
\]
The scoring reference is the better stored return from a dynamic-programming policy and a controller,
$R^\star(s)=\max(R_{\rm DP}(s),R_{\rm ctrl}(s))$.
It is a strong comparator rather than a certificate of optimality.

For policy return $R_\pi(s)$, near-reference and strict-win indicators are
\[
\begin{aligned}
\mathcal N(s)&=\mathbb 1[R_\pi(s)\ge R^\star(s)-5],\\
\mathcal B(s)&=\mathbb 1[R_\pi(s)>R^\star(s)].
\end{aligned}
\]
Task success requires at least 80\% of steps near upright and no not-upright streak longer than 50 steps. A step is near upright when $\cos\theta\ge0.95$ and $|\dot\theta|\le1$.

\paragraph{Evidence set.}
The two-repository census contains 404 standardized evaluations. Every headline row uses the common grid, five independently trained actors, and at most 100,000 selected-route learning transitions. Every displayed scorecard therefore contains 12,505 rollouts. Method and checkpoint selection was retrospective on the standardized campaign, with no independent held-out confirmation. The mixed initializer oversamples a near-upright region in 20\% of its 400,000 static labels. Its later 20,000-transition priority stage selects starts from measured performance.

\paragraph{Technical context and study design.}
SAC trains a stochastic actor and twin critics from reward and entropy \cite{haarnoja2018sac}. SimbaV2 contributes normalized residual networks for stable deep RL \cite{lee2025simba}. DAgger supplies reference labels on learner-visited states to address sequential distribution shift \cite{ross2011dagger}. Our measured-regret priority stage is related to automatic curricula \cite{jiang2021plr}. Published SACn trains on several multi-step losses and uses importance weighting plus a lower-variance entropy estimate to control off-policy bias and target variance \cite{lyskawa2025sacn}. Our FastSACN variant is a simplified ablation: it retains only the one-step and longest losses, and the selected pure recipe does not use the paper's importance correction. Critic action search is a thresholded analogue of critic-guided action optimization \cite{kalashnikov2018qtopt}. Reflection and search are deployment projections rather than new training algorithms. Multi-seed evaluation is essential because aggregate means can conceal instability \cite{henderson2018matters,agarwal2021precipice}.

The study combines these components into two deployments, measures the paired contribution of each retained component, and diagnoses the larger pure-RL tail. The analysis tests four explanations: missing experience, ineffective critic learning, failed actor extraction, and error accumulation across several decisions.

\clearpage
\onecolumn
\small
\section{The two selected recipes}
\vspace{2mm}

\begin{center}
\begin{tikzpicture}[
  node distance=5mm and 4mm,
  every node/.style={font=\small,align=center},
  box/.style={rounded corners=2mm,very thick,minimum height=11mm,text width=25mm},
  mixedbox/.style={box,draw=mixed!70,fill=mixed!7},
  purebox/.style={box,draw=pure!70,fill=pure!7},
  arrow/.style={-{Latex[length=2mm]},very thick,draw=ink!55}
]
\node[mixedbox] (mstatic) {400k static\\reference labels};
\node[mixedbox,right=of mstatic] (mdagger) {30k learner-state\\DAgger transitions};
\node[mixedbox,right=of mdagger] (mpriority) {20k automatic\\priority follow-up};
\node[mixedbox,right=of mpriority] (mactor) {Supervised\\SimbaV2 actor};
\node[mixedbox,right=of mactor] (mlocal) {50k FastSACN8 critic\\and local Q-search};
\draw[arrow] (mstatic) -- (mdagger);
\draw[arrow] (mdagger) -- (mpriority);
\draw[arrow] (mpriority) -- (mactor);
\draw[arrow] (mactor) -- (mlocal);

\node[purebox,below=14mm of mstatic] (psac) {Reward-only SimbaV2\\FastSACN8, 100k};
\node[purebox,right=12mm of psac] (preflect) {Global reflection\\projection};
\node[purebox,right=12mm of preflect] (pq) {41-action\\twin-Q search};
\node[purebox,right=12mm of pq] (paccept) {Both critics must\\improve by 0.005};
\draw[arrow] (psac) -- (preflect);
\draw[arrow] (preflect) -- (pq);
\draw[arrow] (pq) -- (paccept);

\node[font=\bfseries\small,text=mixed!80!black,left=3mm of mstatic] {Mixed};
\node[font=\bfseries\small,text=pure!80!black,left=3mm of psac] {Pure RL};
\end{tikzpicture}
\captionof{figure}{The mixed actor receives direct actions on the states it visits, then a reward critic makes small local corrections. Pure RL learns from reward, imposes the Pendulum symmetry, and searches a wider action set with its critics.}
\end{center}

\vspace{1mm}
\begin{multicols}{2}
\paragraph{Mixed recipe in plain language.}
A width-64, two-block SimbaV2 actor first learns reference actions. Three DAgger rounds then let the actor control the pendulum and label the states it actually reaches. The follow-up evaluates uniformly sampled starts, ranks them by measured regret, and collects 20,000 more transitions from weak starts plus a uniform fraction. The resulting dataset contains states reached after earlier actor errors. Tiny critic-based shifts of the action labels changed none of the requested classifications.

The supervised actor minimizes
\[
\mathcal L_{\rm BC}(\theta)=
\frac{1}{|\mathcal D|}\sum_{(s,a^\star)\in\mathcal D}
\left(\frac{\pi_\theta(s)-a^\star}{2}\right)^2.
\]
At deployment, the actor proposes $a_0=\pi_\theta(s)$. A shared reward-only FastSACN8 critic checks
\[
\mathcal A_{\rm local}=
\operatorname{clip}\!\left(a_0+\{-0.10,-0.05,0,0.05,0.10\},-2,2\right).
\]
The maximizing proposal replaces $a_0$ only when both critics prefer it. The categorical critic loss contains one-step and eight-step cross-entropy terms. Using $w_n=\lambda^{n-1}$ with $\lambda=0.5$, the normalized coefficient on the eight-step term is
\[
\frac{0.5^7}{1+0.5^7}=0.775\%.
\]
This number describes the explicit weighting inside one critic update. It is not a measured gradient share or an effect size. The two cross-entropy terms can have different magnitudes and gradient directions, both update shared critic features, and the update is repeated throughout training. A small explicit coefficient can therefore have a substantial accumulated effect.

\paragraph{Pure-RL recipe in plain language.}
Five independent SimbaV2 actor-critic systems train for 100,000 reward-only transitions. Each actor has width 32 and one residual block. Each twin categorical critic has width 64, two residual blocks, and 51 return atoms \cite{bellemare2017distributional}. Ordinary SAC bootstraps after one reward. Our SACN8 implementation trains on the losses from every horizon from one through eight. FastSACN8 keeps only the one-step and eight-step losses. The selected recipe assigns the eight-step loss coefficient $\lambda^7$ and trains each critic twice per transition. At deployment, symmetry reduces actor inconsistency, then both critics must agree before a searched action replaces the actor proposal.

SAC minimizes $\mathbb E[\alpha\log\pi_\theta(a|s)-\min_iQ_i(s,a)]$ for the actor. For $n\in\{1,8\}$, let $Y_t^{(n)}$ be the $n$-step soft return and let $Z_t^{(n)}$ be its categorical projection:
\[
\begin{aligned}
Y_t^{(n)}={}&\sum_{j=0}^{n-1}\gamma^j
\left(r_{t+j}+\gamma\alpha\widehat H_n(\pi(\cdot|s_{t+j+1}))\right)\\
&+\gamma^n\min_i\bar Q_i(s_{t+n},a'),\qquad a'\sim\pi(\cdot|s_{t+n}).
\end{aligned}
\]
If $\operatorname{CE}(Z,Q_i)$ is the cross-entropy between the projected target distribution and critic $i$, FastSACN8 uses
\[
\mathcal L_Q=
\frac{1}{1+\lambda^7}
\sum_{i=1}^{2}\mathbb E\!\left[
\operatorname{CE}\!\left(Z^{(1)},Q_i\right)+
\lambda^7\operatorname{CE}\!\left(Z^{(8)},Q_i\right)
\right].
\]
Seeds are 0--4. Adam uses batch 256, replay capacity 100k, $\gamma=.99$, $\tau=.005$, and actor/critic rates decaying from $10^{-4}$ to $5{\times}10^{-5}$. The critic receives two updates and the actor one update per environment transition after learning starts.

The deployment projection uses Pendulum sign symmetry:
\[
\begin{aligned}
M(s)&=(\cos\theta,-\sin\theta,-\dot\theta),\\
a_{\rm sym}(s)&=\frac{\pi_\theta(s)-\pi_\theta(M(s))}{2}.
\end{aligned}
\]
The critics then evaluate 41 torques in $[-2,2]$. Proposal
$a_q=\arg\max_a\min_iQ_i(s,a)$ is accepted only if
\[
Q_i(s,a_q)-Q_i(s,a_{\rm sym})>0.005
\quad\text{for }i\in\{1,2\}.
\]
Reflection is one global mathematical rule. Q-search uses learned reward critics and never reads the reference.

\paragraph{Resource accounting.}
The selected mixed route assigns 30k initial DAgger transitions, 20k priority DAgger transitions, and one 50k shared critic, giving 100k learning transitions. Its 690k reference queries comprise 400k initializer labels, 30k learner-state labels, 240k reset-support labels, and 20k priority labels. It also uses 800k candidate-rollout interactions to rank starts. Across five actors, the actual mixed study contains 300k learning transitions because the 50k critic is shared. Pure RL uses 100k learning transitions per seed and 500k across five complete actor-critic runs. The result compares selected pipelines rather than equal total compute.

\begin{center}
\scriptsize
\captionof{table}{Selected-route and full-study resource ledgers.}
\resizebox{\columnwidth}{!}{%
\begin{tabular}{lrrrr}
\toprule
Pipeline & Learn/route & Learn/full & Discover/full & Labels/full\\
\midrule
Auto-priority mixed & 100k & 300k & 4.00m & 3.45m\\
Pure RL & 100k & 500k & 0 & 0\\
\bottomrule
\end{tabular}}
\end{center}
\end{multicols}

\normalsize
\clearpage
\onecolumn
\section{Reliability and component results}

\begin{center}
\includegraphics[width=.55\textwidth]{37_core_scorecard_large.png}
\captionof{figure}{Scorecard on the same 12,505 seed-state trials. Every row uses five actors; the mixed rows share one selected critic.}
\end{center}

\vspace{-2mm}
\begin{minipage}[t]{.46\textwidth}
\centering
\includegraphics[width=\linewidth]{49_failure_footprints_clean.png}
\captionof{figure}{Initial states where at least one of five actors falls more than five return points below the reference.}
\end{minipage}\hfill
\begin{minipage}[t]{.51\textwidth}
\footnotesize
\paragraph{The mixed method controls the reliability tail.}
The selected mixed pipeline reaches 12,496 of 12,505 near-reference trials, compared with 12,008 for pure RL. It records 11,737 task successes versus 11,614. Pure RL records 2,854 strict wins versus 1,570. Reliability and exceeding the stored comparator therefore rank the pipelines differently.

\paragraph{The gap is seed-dependent.}
Mixed near-reference counts range from 2,494 to 2,501 of 2,501 states per seed. Pure counts range from 2,333 to 2,457. The mixed pipeline succeeds for all five seeds on 2,493 cells, or 99.680\%. Pure RL succeeds for all five on 2,293 cells, fails for all five on 13, and varies by seed on the remaining 195 cells.
\end{minipage}

\vspace{1mm}
\begin{minipage}[t]{.44\textwidth}
\centering
\includegraphics[width=\linewidth]{04_component_ablation_report.png}
\captionof{figure}{Matched fixed-minus-broken classification counts.}
\end{minipage}\hfill
\begin{minipage}[t]{.54\textwidth}
\footnotesize
\paragraph{Mixed ablations.}
The 20k learner follow-up adds a net 29 near-reference, seven task-success, and 442 strict-win classifications. Automatic priority adds ten, two, and 111 relative to uniform starts. Local FastSACN8 Q-search adds 26 near-reference and 32 task-success classifications, but loses 243 strict wins. Tiny critic shifts to supervised labels change zero requested classifications.

\paragraph{Pure ablations.}
On the selected FastSACN8 plus SimbaV2 lineage, reflection adds 312 near-reference, ten task-success, and 376 strict-win classifications. Global unanimous Q-search then adds another 458 near-reference, 242 task-success, and 333 strict-win classifications. These paired authority-grid interventions show that the critics and symmetry rule can correct many actions chosen by the actor.
\end{minipage}

\clearpage
\onecolumn
\section{Matched target ablations separate critic direction from actor sensitivity}

\begin{center}
\centering
\includegraphics[width=.98\textwidth]{58_poster_target_family_internals_wide.png}
\captionof{figure}{Matched 100k SimbaV2 target families with five actor-critic seeds each. SAC denotes one-step SAC, Fast denotes FastSACN8, and SACn denotes SACN8. Panels A and B use 5,553 locked states per seed. Panels C and D use 512 locked states per seed and report failure-conditioned interventions. In C, blue is critic-direction agreement with realized return change, green is helpful-reference-action recognition, and orange is the harmful projected-step rate. In D, gray is action-bound occupancy, blue is an outward critic direction, and orange is the fraction of a nominal 0.05 step retained after projection.}
\end{center}

\vspace{-1mm}
\begin{multicols}{2}
\small
\paragraph{Target choice changes critic direction and actor sensitivity differently.}
Relative to one-step SAC, FastSACN8 and SACN8 raise median torque-limit occupancy from 59.1\% to 86.9\% and 86.7\%. Median tanh sensitivity $1-\tanh^2(u)$ falls from 0.120 to 0.039 for both multi-step families. Their critics nevertheless improve local action-direction agreement from 73.1\% to 74.8\% and 76.2\%, while harmful projected steps fall from 27.8\% to 26.6\% and 24.6\%. Better local critic direction therefore coexists with a less responsive actor output.

\paragraph{The action bound absorbs most local corrections on failures.}
Among failure rows, 91.5\% of one-step actions, 98.5\% of FastSACN8 actions, and 98.9\% of SACN8 actions lie within 0.05 torque of the bound. Their critics point farther outward on 82.0\%, 94.6\%, and 95.1\% of those boundary rows. Projection retains 28.2\%, 7.1\%, and 6.9\% of a nominal 0.05 critic-gradient step. This is a post-training, failure-conditioned association. Saturated actions are often correct in Pendulum, so the result identifies a local projection constraint rather than proving that saturation caused failure.

\paragraph{FastSACN8 critic ranking is useful but uneven.}
At the first actor-reference disagreement, one reference action improves continuation return in 79.4\% of FastSACN8 failures, yet the FastSACN8 critic ranks that helpful action above the actor action in 54.5\% of cases. The corresponding recognition rates are 86.4\% for one-step SAC and 80.4\% for SACN8. Reference actions are used only after frozen-policy failures are identified. They never enter deployment.

\paragraph{The selected FastSACN8 gain appears after critic-guided selection.}
On 27,765 locked off-grid trials, five one-step actors reach 91.864\% near-reference success before search, compared with 89.937\% for the selected $\lambda=0.5$, critic-UTD2 FastSACN8 actors. After the same reflection and Q-search rule, the ordering reverses from 94.489\% to 95.909\%. The selected recipe changes both the target and critic update dose, so the comparison does not identify their separate effects. Q-search evaluates candidate torques directly and does not depend on a small projected change to the saturated actor output.

\paragraph{A naive joint imitation and SAC loss is badly scaled.}
In the logged 25,000-step joint pilot, the joint BC plus SACN8 actor reaches 2,443 of 2,501 near-reference starts, compared with 1,757 for the matched reward-only actor. The median norm of the weighted behavior-cloning gradient is 37.2 times the SAC actor-gradient norm. Minibatch cosine similarity ranges from $-0.943$ to $0.935$, with interval median 0.075. The pilot shows that the tested joint loss is severely imbalanced. It does not rule out normalized, scheduled, or projection-based gradient combinations.

\paragraph{Dormancy and replay absence do not explain the observed tail.}
No dormant units are detected at the registered relative-activation threshold. A reference-like action appears among nearby replay actions for 1,010 of 1,021 failures in the established one-step lineage. These negative controls weaken explanations based on widespread dead features or complete absence of corrective actions from replay. They do not show whether rare corrective transitions receive sufficient update weight.

\paragraph{Repair requires a state-matched sequence.}
For each of 1,267 selected raw-actor failures, we use reference actions for the first $K$ decisions and then resume the same frozen actor. Repair rises from 19.1\% at $K=1$ to 47.7\%, 91.3\%, and 99.0\% at $K=8,16,32$. Shuffling the same 32 actions repairs 0.6\%. DAgger records $(s_t,a^\star(s_t))$ at states visited by the current actor, while the priority refit emphasizes starts with measured return deficits. This behavior-level intervention complements the internal diagnostics: local critic information can be useful, yet successful recovery requires a state-dependent action sequence.
\end{multicols}

\clearpage
\twocolumn
